% Options for packages loaded elsewhere
\PassOptionsToPackage{unicode}{hyperref}
\PassOptionsToPackage{hyphens}{url}
\documentclass[
]{article}
\usepackage{xcolor}
\usepackage{amsmath,amssymb}
\setcounter{secnumdepth}{-\maxdimen} % remove section numbering
\usepackage{iftex}
\ifPDFTeX
  \usepackage[T1]{fontenc}
  \usepackage[utf8]{inputenc}
  \usepackage{textcomp} % provide euro and other symbols
\else % if luatex or xetex
  \usepackage{unicode-math} % this also loads fontspec
  \defaultfontfeatures{Scale=MatchLowercase}
  \defaultfontfeatures[\rmfamily]{Ligatures=TeX,Scale=1}
\fi
\usepackage{lmodern}
\ifPDFTeX\else
  % xetex/luatex font selection
\fi
% Use upquote if available, for straight quotes in verbatim environments
\IfFileExists{upquote.sty}{\usepackage{upquote}}{}
\IfFileExists{microtype.sty}{% use microtype if available
  \usepackage[]{microtype}
  \UseMicrotypeSet[protrusion]{basicmath} % disable protrusion for tt fonts
}{}
\makeatletter
\@ifundefined{KOMAClassName}{% if non-KOMA class
  \IfFileExists{parskip.sty}{%
    \usepackage{parskip}
  }{% else
    \setlength{\parindent}{0pt}
    \setlength{\parskip}{6pt plus 2pt minus 1pt}}
}{% if KOMA class
  \KOMAoptions{parskip=half}}
\makeatother
\usepackage{graphicx}
\makeatletter
\newsavebox\pandoc@box
\newcommand*\pandocbounded[1]{% scales image to fit in text height/width
  \sbox\pandoc@box{#1}%
  \Gscale@div\@tempa{\textheight}{\dimexpr\ht\pandoc@box+\dp\pandoc@box\relax}%
  \Gscale@div\@tempb{\linewidth}{\wd\pandoc@box}%
  \ifdim\@tempb\p@<\@tempa\p@\let\@tempa\@tempb\fi% select the smaller of both
  \ifdim\@tempa\p@<\p@\scalebox{\@tempa}{\usebox\pandoc@box}%
  \else\usebox{\pandoc@box}%
  \fi%
}
% Set default figure placement to htbp
\def\fps@figure{htbp}
\makeatother
\setlength{\emergencystretch}{3em} % prevent overfull lines
\providecommand{\tightlist}{%
  \setlength{\itemsep}{0pt}\setlength{\parskip}{0pt}}
\usepackage{bookmark}
\IfFileExists{xurl.sty}{\usepackage{xurl}}{} % add URL line breaks if available
\urlstyle{same}
\hypersetup{
  pdftitle={לוגיקה למדעי המחשב ממן 11},
  hidelinks,
  pdfcreator={LaTeX via pandoc}}

\title{לוגיקה למדעי המחשב ממן 11}
\author{}
\date{}

\begin{document}
\maketitle

לוגיקה למדעי המחשב\\
ממן 11\\
תום ארי

\section{שאלה 1}\label{ux5e9ux5d0ux5dcux5d4-1}

\begin{enumerate}
\tightlist
\item
  המחרוזת {} \textbf{אינה} פסוק מאחר משני סיבות:

  \begin{enumerate}
  \tightlist
  \item
    מספר הסוגריים אינו מאוזן בסיפא של המחרוזת
  \item
    בסיפא של המחרוזת ישנו קשר דו מקומי שאינו מקיים {}. זוהי שגיאה
    תחבירית
  \end{enumerate}
\item
  המחרוזת {} \textbf{אכן} מחרוזת, נראה את עץ הבניה שלה באיור הבא

  \begin{itemize}
  \tightlist
  \item
    \pandocbounded{\includegraphics[keepaspectratio]{/Users/tom/Obsidian Vault/Pasted image 20250323215949.png}}
  \item
    דרגת העץ היא \textbf{6}
  \item
    נבנה סדרות בניה:
  \end{itemize}
\end{enumerate}

\begin{center}\rule{0.5\linewidth}{0.5pt}\end{center}

\begin{enumerate}
\setcounter{enumi}{2}
\tightlist
\item
  המחרוזת {} \textbf{אכן} מחרוזת, נראה את עץ הבנייה שלה באיור הבא

  \begin{itemize}
  \tightlist
  \item
    \pandocbounded{\includegraphics[keepaspectratio]{/Users/tom/Obsidian Vault/Pasted image 20250323220909.png}}
  \item
    דרגת העץ היא \textbf{6}
  \item
    נבנה סדרת בניה
  \end{itemize}
\end{enumerate}

\begin{center}\rule{0.5\linewidth}{0.5pt}\end{center}

\begin{enumerate}
\setcounter{enumi}{3}
\tightlist
\item
  המחרוזת {} \textbf{אינה} מחרוזת משני סיבות:

  \begin{enumerate}
  \tightlist
  \item
    מספר הסוגריים אינו מאוזן
  \item
    בסיפא של המחרוזת ישנו קשר דו מקומי שאינו מקיים {}. זוהי שגיאה
    תחבירית
  \end{enumerate}
\item
  המחרוזת {} \textbf{אינה} פסוק. קל לראות כי {} אינו פסוק, בדומה לדוגמא
  2.1, עמוד 72 בספר לוגיקה למדעי המחשב
\end{enumerate}

\section{שאלה 2}\label{ux5e9ux5d0ux5dcux5d4-2}

\begin{itemize}
\tightlist
\item
  נוכיח בעזרת אינדוקציה מבנית שאם מספר הפסוקים האלמנטריים בפסוק נתון הוא
  {} ומספר הקשרים הדו מקומיים הוא {} אזי {}

  \begin{itemize}
  \tightlist
  \item
    המשוואה {} מתקיימת אםם {} שניהם זוגיים או אי זוגיים בהתאמה. נוכיח
    באינדוקציה על המספרים הטבעיים. עבור {} נקבל {} הנחת האינדוקציה היא
    שעבור פסוקים {} מתקיים {}, {} אתא נפריד למקרים

    \begin{enumerate}
    \tightlist
    \item
      {} פסוק אלמנטרי: לא מתקיים שום סוג קשר ולכן ברור כי {}
    \item
      {} פסוק שלילה: אינו מוסיף פסוק אלמנטרי, אינו קשר דו מקומי. אינו
      משפיע על {}
    \item
      {} קשר דו מקומי: מספר הפסוקים האלמנטריים הוא {}, מספר הקשרים הדו
      מקומיים הוא {} משמע מתקיים {} מאחר והוספנו קשר יחיד. מתוך הנחת
      האינדוקציה נקבל {} משוואה זו נכונה רק אם לשני הצדדים זוגיים או אי
      זוגיים בהתאמה, באופן דומה להגדרה. הראנו כי היחס בין {} מתקיים עבור
      כל מספר פסוקים אלמנטריים, ומספר קשרים דו מקומיים.
    \end{enumerate}
  \end{itemize}
\end{itemize}

\end{document}
